% Suggested addition to Appendix A (after the paragraph ending "Direct prior draws provide an independent check.").
% Not applied to the manuscript.  Numbers: recomputed.json / delta_summary.json of this package.
We tested the sensitivity of the evidences to the event-rate normalisation and to the physical cuts with three separate perturbations, each applied to the three maps of \autoref{tab:single_pointing_evidence}: the H$\alpha$ star-formation rate halved, the same rate doubled, and the cuts widened to $0.01\leq\kappa/({\rm kpc^2\,Gyr^{-1}})\leq5$ and $0.01\leq t_\ast/{\rm Gyr}\leq10$. Only L1 and L2, which sample $\beta_{\rm src}$, were refitted, in every transport interval and with the reference sampler settings and a single seed; L0 and L1n depend on neither quantity and are unchanged. The EB tests reuse the EB covariances of the reference analysis without rebuilding them. Under the null covariance no margin over L0 changes by more than $0.92$ relative to the reference fits (less than $1.5$ across the reference and the three perturbations). For NGC~2466 (S), L2 and L1 lie within $0.9$ of each other and exceed L1n by $3.1$--$4.6$; for NGC~2466 (D16) and NGC~1084, L1n and L2 lie within $0.7$ of each other, in either order for D16, and exceed L1 by more than $5$. Under the EB covariance the L1 and L2 margins change by at most $1.04$ relative to the reference fits; the three extensions of NGC~1084 remain within $1.4$ of one another and are not distinguished, and no margin for NGC~2466 (D16) exceeds $1.6$. These results apply only to the three perturbations tested, with per-run sampler uncertainties of $0.12$--$0.24$ in $\ln{\cal Z}$, and do not establish that the evidences are independent of the prior ranges in general.
